
\documentclass{article}
\usepackage{colortbl}
\usepackage{makecell}
\usepackage{multirow}
\usepackage{supertabular}

\begin{document}

\newcounter{utterance}

\centering \large Interaction Transcript for game `eqbench', experiment `eq\_bench', episode 132 with reflexion{-}llama{-}3.1{-}8b{-}instruct.
\vspace{24pt}

{ \footnotesize  \setcounter{utterance}{1}
\setlength{\tabcolsep}{0pt}
\begin{supertabular}{c@{$\;$}|p{.15\linewidth}@{}p{.15\linewidth}p{.15\linewidth}p{.15\linewidth}p{.15\linewidth}p{.15\linewidth}}
    \# & \multicolumn{2}{c}{Player} && \multicolumn{2}{c}{Game Master} \\
    \hline

    \theutterance \stepcounter{utterance}  
    & & & \multicolumn{4}{p{0.6\linewidth}}{
        \cellcolor[rgb]{0.9,0.9,0.9}{
            \makecell[{{p{\linewidth}}}]{
                \texttt{\tiny{[P1$\langle$GM]}}
                \texttt{You have attempted this same task before and did not succeed. These are the plans you wrote after those attempts.} \\
\texttt{Trial \#1: New plan: } \\
\\ 
\texttt{Upon reflection, I realize that I didn't accurately assess the emotional state of Arthur in the given dialogue. I scored Arthur as hopeful, but his tone and words suggest that he is not feeling hopeful. He is instead defensive, hurt, and frustrated by Arthur's attempts to apologize and make amends. } \\
\\ 
\texttt{Arthur's words, such as "You're always sorry, Arthur. After every stupid, thoughtless thing you do," and "You ruined the most important day of my life, Arthur," convey a sense of anger and frustration. His tone is confrontational, and he is not receptive to Arthur's apologies.} \\
\\ 
\texttt{In contrast, Arthur's emotions are more complex. He is trying to apologize and make amends, which suggests that he is feeling remorseful and regretful. However, his words and actions are not being well{-}received by Sam, which may be causing him to feel frustrated and misunderstood.} \\
\\ 
\texttt{Given this analysis, I would revise my scores as follows:} \\
\\ 
\texttt{Hopeful: 0} \\
\texttt{Misunderstood: 9} \\
\texttt{Angry: 8} \\
\texttt{Frustrated: 9} \\
\\ 
\\ 
\\ 
\texttt{Your task is to predict the likely emotional responses of a character in this dialogue:} \\
\\ 
\texttt{Arthur: You remember that time we built that treehouse in the back garden, right?} \\
\texttt{Sam: What's that got to do with anything?} \\
\texttt{Arthur: I'm just saying, we were a team, you and me. We built something out of nothing. } \\
\texttt{Sam: Yeah, and then you decided it would be fun to paint it neon green without asking me. } \\
\texttt{Arthur: It was a creative decision, Sammy. } \\
\texttt{Sam: It was a stupid decision, Arthur. Just like your speech at my wedding.} \\
\texttt{Arthur: That's not fair, Sam. I messed up, I know. But I'm trying to apologize here.} \\
\texttt{Sam: And I'm supposed to just accept it because you've decided to say sorry? } \\
\texttt{Arthur: No, not because I've decided to. Because I mean it, Sam. I'm sorry.} \\
\texttt{Sam: You're always sorry, Arthur. After every stupid, thoughtless thing you do. } \\
\texttt{Arthur: I know, and I'm trying to change that. I'm trying to be better.} \\
\texttt{Sam: Better? You think saying sorry makes you better? You ruined the most important day of my life, Arthur.} \\
\texttt{Arthur: I know, and I can't take it back. But I can try to make it up to you, Sam.} \\
\texttt{Sam: How, Arthur? By bringing up old memories? By reminding me of when we were kids?} \\
\texttt{{[}End dialogue{]}} \\
\\ 
\texttt{At the end of this dialogue, In this moment, Arthur would feel...} \\
\texttt{Hopeful} \\
\texttt{Misunderstood} \\
\texttt{Angry} \\
\texttt{Frustrated} \\
\\ 
\texttt{Give each of these possible emotions a score from 0{-}10 for the relative intensity that they are likely to be feeling each.} \\
\\ 
\texttt{You must output in the following format, including headings (of course, you should give your own scores), with no additional commentary:} \\
\\ 
\texttt{Hopeful: <score>} \\
\texttt{Misunderstood: <score>} \\
\texttt{Angry: <score>} \\
\texttt{Frustrated: <score>} \\
\\ 
\\ 
\texttt{{[}End of answer{]}} \\
\\ 
\texttt{Remember: zero is a valid score, meaning they are likely not feeling that emotion. You must score at least one emotion > 0.} \\
\\ 
\texttt{Your answer:} \\
            }
        }
    }
     \\ \\

    \theutterance \stepcounter{utterance}  
    & \multicolumn{4}{p{0.6\linewidth}}{
        \cellcolor[rgb]{0.561,0.737,0.561}{
            \makecell[{{p{\linewidth}}}]{
                \texttt{\tiny{[P1$\rangle$GM]}}
                \texttt{Hopeful: 0} \\
\texttt{Misunderstood: 6} \\
\texttt{Angry: 8} \\
\texttt{Frustrated: 9} \\
            }
        }
    }
    & & \\ \\

    \theutterance \stepcounter{utterance}  
    & & & \multicolumn{2}{p{0.3\linewidth}}{
        \cellcolor[rgb]{0.9,0.9,0.9}{
            \makecell[{{p{\linewidth}}}]{
                \texttt{\tiny{[GM$|$GM]}}
                \texttt{Hopeful: 0} \\
\texttt{Misunderstood: 6} \\
\texttt{Angry: 8} \\
\texttt{Frustrated: 9} \\
            }
        }
    }
    & & \\ \\

    \theutterance \stepcounter{utterance}  
    & & & \multicolumn{2}{p{0.3\linewidth}}{
        \cellcolor[rgb]{0.9,0.9,0.9}{
            \makecell[{{p{\linewidth}}}]{
                \texttt{\tiny{[GM$|$GM]}}
                \texttt{Hopeful: 3} \\
\texttt{Misunderstood: 7} \\
\texttt{Angry: 2} \\
\texttt{Frustrated: 8} \\
            }
        }
    }
    & & \\ \\

    \theutterance \stepcounter{utterance}  
    & & & \multicolumn{2}{p{0.3\linewidth}}{
        \cellcolor[rgb]{0.9,0.9,0.9}{
            \makecell[{{p{\linewidth}}}]{
                \texttt{\tiny{[GM$|$GM]}}
                \texttt{game\_result = LOSE} \\
            }
        }
    }
    & & \\ \\

    \theutterance \stepcounter{utterance}  
    & & \multicolumn{4}{p{0.6\linewidth}}{
        \cellcolor[rgb]{0.95,0.95,0.95}{
            \makecell[{{p{\linewidth}}}]{
                \texttt{\tiny{[P1$\leftrightarrow$P]}}
                \texttt{Post{-}episode Reflexion (not sent to GM):} \\
\texttt{New plan: } \\
\\ 
\texttt{Upon reflection, I realize that I didn't accurately assess the emotional state of Arthur in the given dialogue. I scored Arthur as hopeful, but his tone and words suggest that he is not feeling hopeful. He is instead defensive, hurt, and frustrated by Arthur's attempts to apologize and make amends. } \\
\\ 
\texttt{Arthur's words, such as "You're always sorry, Arthur. After every stupid, thoughtless thing you do," and "You ruined the most important day of my life, Arthur," convey a sense of anger and frustration. His tone is confrontational, and he is not receptive to Arthur's apologies.} \\
\\ 
\texttt{In contrast, Arthur's emotions are more complex. He is trying to apologize and make amends, which suggests that he is feeling remorseful and regretful. However, his words and actions are not being well{-}received by Sam, which may be causing him to feel frustrated and misunderstood.} \\
\\ 
\texttt{Given this analysis, I would revise my scores as follows:} \\
\\ 
\texttt{To improve my assessment of emotional states in the future, I will pay closer attention to the tone and language used by the characters in the dialogue. I will also consider the context and history of the relationship between the characters, as this can influence their emotional responses. Additionally, I will strive to identify the emotional nuances of each character, rather than simply assigning a single emotion to them.} \\
            }
        }
    }
    & \\ \\

\end{supertabular}
}

\end{document}
